\subsubsection{QEMU Object Model (QOM)}
\label{chap:QEMU_QOM}

The QEMU Object Model (QOM) subsystem is QEMU's object-oriented framework. 
This subsystem provides a device modeling abstraction layer by implements a 
type system using pure C which, despite C's lack of native object-oriented programming 
(OOP) features, allows dynamic type registration, inheritance, and runtime 
polymorphism, providing\cite{Bellard2005Qemu} \cite{Qemu_QOM}:

\begin{itemize}
    \item \textbf{Object-oriented programming:    }Without requiring an OOP 
                                                    language like C++.

    \item \textbf{Dynamic type registration:      }Types can be registered at runtime without 
                                                    recompilation.
    
    \item \textbf{Inheritance and polymorphism:   }Types inherit from parent types (single 
                                                    inheritance) and override methods.
    
    \item \textbf{Runtime reflection:             }All objects and properties are 
                                                    discoverable via standardized queries.
    
    \item \textbf{Dynamic properties:             }Each object instance has configurable, 
                                                    typed properties.
    
    \item \textbf{Composition tree:               }Hierarchical object relationships enabling complex 
                                                    system modeling.
\end{itemize}

Before the introduction of QOM, QEMU development was fragmented. Each 
board implemented its devices independently, requiring duplicate code; device 
configurations were hardcoded; and extending QEMU with new hardware was cumbersome. 
The introduction of the QOM subsystem resolved these problems by unifying 
device management through the implementation of:

\begin{itemize}
    \item \textbf{Uniform object model:             }All devices are Objects with properties.
    
    \item \textbf{Runtime type registration:        }Types defined via TypeInfo at startup.
    
    \item \textbf{Common instantiation interface:   }Devices created uniformly by type name.
    
    \item \textbf{Dynamic configuration:            }Properties accessible and modifiable via a 
                                                        standardized API.
\end{itemize}


QOM's type system is foundational to QEMU device modeling. Every 
emulated peripheral (timers, ADCs, I2C controllers) is a QOM Object instance with 
an ObjectClass defining its behavior. This abstraction decouples device logic from 
QEMU core, enabling researchers to develop and validate new device models—such 
as the STM32F405 I2C controller developed in this thesis—within a consistent, 
well-defined framework. The figure~\ref{fig:qom_management} illustrates the scope of QOM's management 
within QEMU.

\begin{figure}[H]
    \centering
%    \includegraphics[width=0.7\textwidth]{Resources/Draw/Full-system.pdf}
    \caption{QOM Management Scope}
    \label{fig:qom_management}
\end{figure}

There are several key elements that need to be known and understood about 
the QOM API, such as the \texttt{ObjectClass}, \texttt{Object}, \texttt{ObjectProperty} and \texttt{TypeInfo} 
structures, which were used during the implementation and development of this 
thesis. These and other elements, which will be mentioned throughout the thesis, 
are explained in the Appendix~\ref{ann:qemu_c_structures}.
